% 20260526_assingment_Algebra_IIA
\documentclass[dvipdfmx]{jsreport}
\usepackage{soupreamble}
\title{代数学特論IIA 第1回レポート課題}
\author{M1 6012 川村聡一郎}
\date{}

\begin{document}
\maketitle

\begin{tcolorbox}[title= レポート課題1]
	$\mu$を正実数とする。このとき次を示せ。
\begin{enumerate}
	\item $\sum_{k= 1}^\infty \frac{1}{k^\mu}$は$\mu> 1$のとき収束する。
	\item $\sum_{k= 1}^\infty \frac{1}{k^\mu}$は$\mu\le 1$のとき発散する。
\end{enumerate}
\end{tcolorbox}

\begin{souanswer}
\begin{enumerate}
	\item 関数$f(x)= \frac{1}{x^\mu}$を考える。$\mu> 0$であるから、区間$[1, \infty)$で$f(z)$は正かつ連続で単調減少．
\begin{equation*}
	\int_1^M f(x)dx= \int_1^M \frac{1}{x^\mu} dx= \left[\frac{x^{1- \mu}}{1- \mu}\right]_1^M= \frac{M^{1- \mu}}{1- \mu}
\end{equation*}
	ここで$1- \mu< 0\ (\because \mu> 1)$。したがって$M\to \infty$で$M^{1- \mu}\to 0$。よって$\sum_{k= 1}^\infty \frac{1}{k^\mu}$は$\mu> 1$で収束。
	
	\item 
\begin{align*}
	\sum_{k= 1}^\infty \frac{1}{k^\mu}
	&= \frac{1}{1^\mu}+ \frac{1}{2^\mu}+ \frac{1}{3^\mu}+ \frac{1}{4^\mu}+ \cdots\\
	&> \frac{1}{1^\mu}+ \frac{1}{2^\mu}+ \frac{1}{2^\mu}+ \frac{1}{4^\mu}+ \cdots\\
	&= \frac{1}{2^{0\cdot \mu}}+ \frac{1}{2^{1\cdot \mu}}\times 2+ \frac{1}{2^{2\cdot \mu}}\times 4+ \cdots\\
\end{align*}
	これは初項1、公比$2^{1- \mu}(> 0)$の等比級数であり、この級数の和は$\frac{1\times\left(1- \left(\frac{2}{2^\mu}\right)^n\right)}{1- \frac{2}{2^\mu}}\to \infty\ (n\to \infty)$となるから、$\sum_{k= 1}^\infty \frac{1}{k^\mu}$は$\mu\le 1$で発散する。
\end{enumerate}	
\end{souanswer}


\begin{tcolorbox}[title= レポート課題2]
	$\omega_1, \omega_2\in \C$に対し、次が成立することを示せ。
\begin{enumerate}
	\item $\omega_1\ne 0\Longleftrightarrow \omega_1$は$\C$上一次独立
	\item $\omega_1$と$\omega_2$が$\R$上一次独立$\Longleftrightarrow \frac{\omega_1}{\omega_2}\not\in \R$以下が同値であることを示せ。
\end{enumerate}
\end{tcolorbox}

\begin{souanswer} %答案はここに書く
	答案はここに書く
\begin{enumerate}
	\item 
\end{enumerate}	
\end{souanswer}









\end{document}